\documentclass[11pt]{article}

% Shared notes settings for Elements of AIML lectures (UPES).
% Keep this file free of \documentclass to allow reuse via \input.

\usepackage[utf8]{inputenc}
\usepackage[T1]{fontenc}
\usepackage[a4paper,margin=1in]{geometry}
\usepackage{hyperref}
\usepackage{xurl}
\usepackage{booktabs}
\usepackage{amsmath}
\usepackage{amssymb}
\usepackage{listings}
\usepackage{xcolor}
\usepackage{enumitem}
\usepackage{parskip}

% Code listings (general-purpose).
\lstset{
  basicstyle=\ttfamily\small,
  keywordstyle=\color{blue},
  commentstyle=\color{gray},
  stringstyle=\color{teal},
  showstringspaces=false,
  columns=fullflexible,
  frame=single,
  framerule=0.2pt,
  breaklines=true
}

\setlist[itemize]{topsep=0.3em,itemsep=0.2em}
\setlist[enumerate]{topsep=0.3em,itemsep=0.2em}

% Simple per-section source line.
\newcommand{\notesources}[1]{\par\noindent{\small\color{gray}\textit{Sources:} \sloppy #1\par}}

% Unit-wise source strings for this subject (used by slides/notes).
% Path is relative to the lecture `latex/` folder (the compilation CWD).
\IfFileExists{../../../_shared/aiml-sources.tex}{%
  % Source strings for Elements of AIML (used by slides + notes).

\newcommand{\AIMLRepoURL}{https://github.com/tali7c/Elements-of-AIML}

\newcommand{\AIMLLabSources}{%
Provost and Fawcett, \textit{Data Science for Business}.;%
McKinney, \textit{Python for Data Analysis}.;%
WEKA Documentation (Waikato Environment for Knowledge Analysis), accessed 2026-02-28.;%
Matplotlib and Seaborn documentation, accessed 2026-02-28.%
}

\newcommand{\AIMLUnitISources}{%
Rich and Knight, \textit{Artificial Intelligence}, McGraw Hill, 2017.;%
Russell and Norvig, \textit{Artificial Intelligence: A Modern Approach} (4e), Ch. 1--2 (Introduction, Intelligent Agents).;%
Poole and Mackworth, \textit{Artificial Intelligence: Foundations of Computational Agents} (2e), selected sections.%
}

\newcommand{\AIMLUnitIISources}{%
Russell and Norvig, \textit{Artificial Intelligence: A Modern Approach} (4e), Ch. 7--9 (Logical Agents, First-Order Logic, Inference).;%
Huth and Ryan, \textit{Logic in Computer Science} (2e), selected sections on propositional and first-order logic.;%
Genesereth and Nilsson, \textit{Logical Foundations of Artificial Intelligence}, selected chapters.%
}

\newcommand{\AIMLUnitIIISources}{%
Mueller and Massaron, \textit{Machine Learning for Dummies}, For Dummies, 2016.;%
Mitchell, \textit{Machine Learning}, selected chapters on learning basics and evaluation.;%
Scikit-learn documentation, accessed 2026-02-28.%
}

\newcommand{\AIMLUnitIVSources}{%
Mueller and Massaron, \textit{Machine Learning for Dummies}, For Dummies, 2016.;%
James, Witten, Hastie, and Tibshirani, \textit{An Introduction to Statistical Learning} (2e), selected chapters.;%
Scikit-learn documentation, accessed 2026-02-28.%
}

\newcommand{\AIMLUnitVSources}{%
NIST, \textit{AI Risk Management Framework (AI RMF 1.0)}, 2023.;%
Barocas, Hardt, and Narayanan, \textit{Fairness and Machine Learning} (online book), selected sections.;%
Knaflic, \textit{Storytelling with Data}, selected chapters on communicating results.%
}
%
}{%
  % Faculty copies live under `private/` so their compile CWD is different.
  \IfFileExists{../../../../public/_shared/aiml-sources.tex}{%
    % Source strings for Elements of AIML (used by slides + notes).

\newcommand{\AIMLRepoURL}{https://github.com/tali7c/Elements-of-AIML}

\newcommand{\AIMLLabSources}{%
Provost and Fawcett, \textit{Data Science for Business}.;%
McKinney, \textit{Python for Data Analysis}.;%
WEKA Documentation (Waikato Environment for Knowledge Analysis), accessed 2026-02-28.;%
Matplotlib and Seaborn documentation, accessed 2026-02-28.%
}

\newcommand{\AIMLUnitISources}{%
Rich and Knight, \textit{Artificial Intelligence}, McGraw Hill, 2017.;%
Russell and Norvig, \textit{Artificial Intelligence: A Modern Approach} (4e), Ch. 1--2 (Introduction, Intelligent Agents).;%
Poole and Mackworth, \textit{Artificial Intelligence: Foundations of Computational Agents} (2e), selected sections.%
}

\newcommand{\AIMLUnitIISources}{%
Russell and Norvig, \textit{Artificial Intelligence: A Modern Approach} (4e), Ch. 7--9 (Logical Agents, First-Order Logic, Inference).;%
Huth and Ryan, \textit{Logic in Computer Science} (2e), selected sections on propositional and first-order logic.;%
Genesereth and Nilsson, \textit{Logical Foundations of Artificial Intelligence}, selected chapters.%
}

\newcommand{\AIMLUnitIIISources}{%
Mueller and Massaron, \textit{Machine Learning for Dummies}, For Dummies, 2016.;%
Mitchell, \textit{Machine Learning}, selected chapters on learning basics and evaluation.;%
Scikit-learn documentation, accessed 2026-02-28.%
}

\newcommand{\AIMLUnitIVSources}{%
Mueller and Massaron, \textit{Machine Learning for Dummies}, For Dummies, 2016.;%
James, Witten, Hastie, and Tibshirani, \textit{An Introduction to Statistical Learning} (2e), selected chapters.;%
Scikit-learn documentation, accessed 2026-02-28.%
}

\newcommand{\AIMLUnitVSources}{%
NIST, \textit{AI Risk Management Framework (AI RMF 1.0)}, 2023.;%
Barocas, Hardt, and Narayanan, \textit{Fairness and Machine Learning} (online book), selected sections.;%
Knaflic, \textit{Storytelling with Data}, selected chapters on communicating results.%
}
%
  }{%
    % Source strings for Elements of AIML (used by slides + notes).

\newcommand{\AIMLRepoURL}{https://github.com/tali7c/Elements-of-AIML}

\newcommand{\AIMLLabSources}{%
Provost and Fawcett, \textit{Data Science for Business}.;%
McKinney, \textit{Python for Data Analysis}.;%
WEKA Documentation (Waikato Environment for Knowledge Analysis), accessed 2026-02-28.;%
Matplotlib and Seaborn documentation, accessed 2026-02-28.%
}

\newcommand{\AIMLUnitISources}{%
Rich and Knight, \textit{Artificial Intelligence}, McGraw Hill, 2017.;%
Russell and Norvig, \textit{Artificial Intelligence: A Modern Approach} (4e), Ch. 1--2 (Introduction, Intelligent Agents).;%
Poole and Mackworth, \textit{Artificial Intelligence: Foundations of Computational Agents} (2e), selected sections.%
}

\newcommand{\AIMLUnitIISources}{%
Russell and Norvig, \textit{Artificial Intelligence: A Modern Approach} (4e), Ch. 7--9 (Logical Agents, First-Order Logic, Inference).;%
Huth and Ryan, \textit{Logic in Computer Science} (2e), selected sections on propositional and first-order logic.;%
Genesereth and Nilsson, \textit{Logical Foundations of Artificial Intelligence}, selected chapters.%
}

\newcommand{\AIMLUnitIIISources}{%
Mueller and Massaron, \textit{Machine Learning for Dummies}, For Dummies, 2016.;%
Mitchell, \textit{Machine Learning}, selected chapters on learning basics and evaluation.;%
Scikit-learn documentation, accessed 2026-02-28.%
}

\newcommand{\AIMLUnitIVSources}{%
Mueller and Massaron, \textit{Machine Learning for Dummies}, For Dummies, 2016.;%
James, Witten, Hastie, and Tibshirani, \textit{An Introduction to Statistical Learning} (2e), selected chapters.;%
Scikit-learn documentation, accessed 2026-02-28.%
}

\newcommand{\AIMLUnitVSources}{%
NIST, \textit{AI Risk Management Framework (AI RMF 1.0)}, 2023.;%
Barocas, Hardt, and Narayanan, \textit{Fairness and Machine Learning} (online book), selected sections.;%
Knaflic, \textit{Storytelling with Data}, selected chapters on communicating results.%
}
%
  }%
}


\title{Elements of AIML Lab\\Experiment 07 (After-submission Reference): Scaling and One-Hot Encoding}
\author{Tofik Ali}
\date{\today}

\begin{document}
\maketitle
\tableofcontents

\section{How to use this reference}
This reference is meant to be shared \textbf{after you submit} Experiment~07.
It explains the two most common preprocessing operations in ML:
\begin{itemize}[leftmargin=*]
  \item scaling numeric features, and
  \item one-hot encoding categorical features,
\end{itemize}
while highlighting the most important risk: \textbf{data leakage}.

\section{Scaling (what it does)}
\subsection{Standardization (z-score)}
Transforms a numeric feature $x$ to:
\[
  z = \frac{x - \mu}{\sigma}
\]
so the transformed feature has mean $\approx 0$ and standard deviation $\approx 1$ on the training data.

\subsection{Min-max normalization}
Transforms to:
\[
  x' = \frac{x - \min(x)}{\max(x) - \min(x)}
\]
so values map to $[0,1]$ (on the training data range).

\section{One-hot encoding (what it does)}
A categorical column like \texttt{City} with values \texttt{\{Delhi, Mumbai\}} becomes two binary columns such as \texttt{City\_Delhi}, \texttt{City\_Mumbai}.

\section{Worked example (leakage-safe pipeline)}
\subsection{Step 1: load and split}
\begin{lstlisting}[language=Python]
import pandas as pd
from sklearn.model_selection import train_test_split

df = pd.read_csv("processed.csv")  # or your dataset

X = df.drop(columns=["target"])    # change target name
y = df["target"]

X_train, X_test, y_train, y_test = train_test_split(
    X, y, test_size=0.2, random_state=42, stratify=y
)
\end{lstlisting}

\subsection{Step 2: build a ColumnTransformer + Pipeline}
\begin{lstlisting}[language=Python]
from sklearn.compose import ColumnTransformer
from sklearn.preprocessing import OneHotEncoder, StandardScaler
from sklearn.pipeline import Pipeline
from sklearn.linear_model import LogisticRegression

num_cols = X_train.select_dtypes("number").columns
cat_cols = X_train.select_dtypes(exclude="number").columns

pre = ColumnTransformer([
    ("num", StandardScaler(), num_cols),
    ("cat", OneHotEncoder(handle_unknown="ignore"), cat_cols),
])

pipe = Pipeline([
    ("pre", pre),
    ("clf", LogisticRegression(max_iter=1000)),
])

pipe.fit(X_train, y_train)
\end{lstlisting}
\textbf{Why this is correct:} scaling and encoding are learned from \textbf{training} and applied to \textbf{test} without leaking test statistics.

\subsection{Step 3: inspect the transformed feature space (optional but useful)}
\begin{lstlisting}[language=Python]
Xtr = pipe.named_steps["pre"].transform(X_train)
print("Transformed train shape:", Xtr.shape)
\end{lstlisting}
If you need column names:
\begin{lstlisting}[language=Python]
feat_names = pipe.named_steps["pre"].get_feature_names_out()
print(feat_names[:20])
\end{lstlisting}

\section{What to write in reports (good explanations)}
\begin{itemize}[leftmargin=*]
  \item ``We standardized numeric features to remove scale effects, because our model uses gradients/distances.''
  \item ``We used one-hot encoding with \texttt{handle\_unknown='ignore'} so unseen categories in test do not crash the model.''
  \item ``We used a pipeline to avoid leakage: preprocessing parameters were fit using training data only.''
\end{itemize}

\section{Troubleshooting}
\begin{itemize}[leftmargin=*]
  \item \textbf{Pipeline error:} check column names and ensure \texttt{target} exists.
  \item \textbf{Too many one-hot columns:} remove ID-like categorical features; group rare categories.
  \item \textbf{Model performs worse after scaling:} some models (trees) do not need scaling; explain model choice.
\end{itemize}

\notesources{\AIMLLabSources}
\end{document}

